\section{Design}
\label{sec:design}

We plan and report the study with the ADEMP structure of \citet{morris2019simulation}: aims, data-generating mechanisms, estimands, methods and performance measures.

\subsection{Aims}
For each validator and each return family the study estimates (i) the probability that the validator calls the best of $N$ skill-less strategies skilled (its \emph{size}), and (ii) the probability that it rejects when one of the $N$ strategies has real skill (its \emph{power}). A validator is \emph{calibrated} in a family when its size is within Monte Carlo error of its nominal level.

\subsection{Data-generating mechanisms}
One replication draws a \emph{search}: $N=\nz{trials}$ return series of $T=\nz{obs}$ periods, read as two years of daily returns ($q=252$ periods per year). Every series has population mean zero and population variance one, so every strategy's true Sharpe ratio is zero. In the skill arms, series 0 receives a constant drift of $s/\sqrt{252}$ per period, which gives it a true annualized Sharpe ratio of $s\in\{1,2,3\}$; the other series are unchanged. The nine families are:

\begin{description}
  \item[IID normal.] Independent standard normal returns; the case every correction is derived for.
  \item[Student $t_4$.] Student $t$ with 4 degrees of freedom, scaled to unit variance: symmetric fat tails. Its fourth moment is infinite, so sample kurtosis does not converge.
  \item[Skew $\pm1.32$.] A centred, unit-variance lognormal with log-standard deviation 0.4, whose skewness is $+1.32$, and its mirror image ($-1.32$).
  \item[GARCH(1,1).] $h_t = 0.05 + 0.10\,\varepsilon_{t-1}^2 + 0.85\,h_{t-1}$, $\varepsilon_t=\sqrt{h_t}\,z_t$ \citep{bollerslev1986garch}, with unconditional variance one; each series starts at $h=1$, $\varepsilon=0$, without burn-in.
  \item[AR(1).] $x_t = 0.2\,x_{t-1} + \sqrt{1-0.2^2}\,z_t$, started from its stationary distribution.
  \item[Two volatility regimes.] Normal returns whose standard deviation is 0.50 or 1.32 (0.6 and 1.6 rescaled to unit variance), switching with probability 0.02 each period.
  \item[Correlated trials.] $x_{k,t} = \sqrt{1/2}\,f_t + \sqrt{1/2}\,z_{k,t}$ with a factor $f_t$ common to all strategies, so every pair of strategies has correlation 0.5.
  \item[Block clusters.] Four independent clusters of five strategies; strategies in a cluster share a factor with pairwise correlation 0.8. Twenty strategies that are really about four ideas, the shape a parameter grid around a few rules produces.
\end{description}

Instrument tests run before any validator is scored: every family's sample mean, variance, skewness, lag-one autocorrelation, autocorrelation of squared returns and cross-strategy correlation are checked against the property it claims, and the drift is checked to give series 0, and only series 0, the intended Sharpe ratio.

\subsection{Estimands}
In each replication the validator sees all \nz{trials} series and returns a $p$-value for the hypothesis that the best strategy has no skill; the single-series validators select the strategy with the highest sample Sharpe ratio, the joint tests consider all strategies at once. The estimands are the rejection rates $\Pr(p \le 0.05)$ in the null arm (size) and the skill arms (power). Because raw power rewards a liberal test, we also report size-adjusted power: the power at the cutoff at which the same validator rejects at most 5\% of the matching null cell. A user cannot apply that cutoff without knowing the family, but it shows how well each statistic ranks skilled searches above skill-less ones.

\subsection{Methods}
Let $\widehat{SR}$ be the best strategy's per-period Sharpe ratio (sample standard deviation with divisor $T-1$), $\widehat{SR}_a=\widehat{SR}\sqrt{q}$ its annualized value, and $\hat\gamma_3,\hat\gamma_4$ its sample skewness and (non-excess) kurtosis. The single-series validators are:

\begin{description}
  \item[LET.] The luck-equivalent-trials test. Under the null and normal returns, $t=\widehat{SR}\sqrt{T}$ is Student $t$ with $T-1$ degrees of freedom, so one skill-less strategy reaches the observed value with probability $p_1=\Pr(t_{T-1}\ge t)$, and the best of $N$ independent ones with probability $1-(1-p_1)^N$ \citep{sidak1967}, the reported $p$-value. As a test it is the \v{S}id\'ak-adjusted Student $t$ test; the statistic reports the same evidence in units of trials, $N_q=\ln(1-q)/\ln(1-p_1)$.
  \item[LET-Lo.] LET after deflating $\widehat{SR}_a$ by the autocorrelation-adjusted annualisation factor of \citet{lo2002sharpe}, $\big[1+\tfrac{2\hat\rho}{1-\hat\rho}\big(1-\tfrac{1-\hat\rho^{\,q}}{q(1-\hat\rho)}\big)\big]^{-1/2}$, at the best series' lag-one autocorrelation $\hat\rho$ (clamped to $[-0.95,0.95]$), as the haircut code of \citet{harvey2015backtesting} does. The factor assumes $\rho_k=\rho^k$, which the AR(1) family satisfies exactly.
  \item[LET-NN.] LET with the normal approximation and the non-normal standard error of \citet{mertens2002comments} evaluated at the estimate: $z=\widehat{SR}\sqrt{T-1}\,/\sqrt{1-\hat\gamma_3\widehat{SR}+\tfrac{\hat\gamma_4-1}{4}\widehat{SR}^2}$.
  \item[HC, HC-Lo, Holm.] The haircut Sharpe ratio of \citet{harvey2015backtesting}: the best strategy's two-sided Student $t$ $p$-value, adjusted by Bonferroni for \nz{trials} tests, without and with the Lo factor, and the Holm \citep{holm1979} variant. Harvey and Liu define the haircut with two-sided $p$-values, which for a positive Sharpe ratio is a one-sided test at 2.5\%; to compare tests at a common level we also read it one-sided (two-sided $p$ halved), and the tables report that reading unless stated.
  \item[DSR.] The deflated Sharpe ratio \citep{bailey2014dsr}: the probabilistic Sharpe ratio \citep{bailey2012frontier} against a benchmark equal to the expected maximum Sharpe ratio of $N$ skill-less trials, computed with the best strategy's skewness and kurtosis and the cross-trial standard deviation of the \nz{trials} annualized Sharpe ratios. Its authors accept a strategy when $\mathrm{DSR}\ge 0.95$; we score $1-\mathrm{DSR}$ as its $p$-value, so that rejection at 5\% is exactly that rule.
  \item[Boot.] An i.i.d.\ bootstrap \citep{efron1979bootstrap} of the best strategy's demeaned returns with $B=\nz{v1.boot_best_draws}$ resamples; the single-trial $p$-value $(h+1)/(B+1)$, where $h$ counts resampled Sharpe ratios at or above the observed one, is \v{S}id\'ak-adjusted for \nz{trials} trials.
  \item[HAC-$t$.] The best strategy's mean divided by a Newey--West (Bartlett kernel) standard error \citep{newey1987simple} with lag $\lfloor 4(T/100)^{2/9}\rfloor$, read against Student $t$ with $T-1$ degrees of freedom and \v{S}id\'ak-adjusted.
  \item[Boot-$t$.] A stationary-bootstrap $t$ test of the best strategy's mean \citep{politis1994stationary} with mean block length $\mathrm{round}(T^{1/3})$ and $\nz{v1.boot_t_reps}$ resamples, each resample studentized by its own standard deviation; \v{S}id\'ak-adjusted.
\end{description}

The joint resampling tests use all \nz{trials} series and one stationary bootstrap of the whole search (the same resampled periods for every strategy, so the dependence between strategies is kept), with mean block length $b=\mathrm{round}(T^{1/3})=8$ and \nz{v1.spa_reps} resamples. With $\bar d_k$ the mean of strategy $k$ and $\bar d_k^*$ its mean in a resample:
\begin{description}
  \item[RC.] White's Reality Check \citep{white2000reality}: $\max_k \bar d_k$ against the resampled distribution of $\max_k(\bar d_k^*-\bar d_k)$. It is not studentized.
  \item[SPA$_c$, SPA$_u$.] Hansen's test for superior predictive ability \citep{hansen2005spa}: $T^{\mathrm{SPA}}=\max\big(0,\max_k \sqrt{T}\bar d_k/\hat\omega_k\big)$, where $\hat\omega_k^2=\hat\gamma_{0,k}+2\sum_{i=1}^{T-1}\kappa_i\hat\gamma_{i,k}$, $\kappa_i=(1-\tfrac iT)(1-\tfrac1b)^i+\tfrac iT(1-\tfrac1b)^{T-i}$, is the Politis--Romano estimate of the long-run variance at the bootstrap's block length; the resampled statistic is $\max\big(0,\max_k\sqrt{T}(\bar d_k^*-g_k)/\hat\omega_k\big)$ with the same $\hat\omega_k$. The consistent $p$-value recentres at $g_k=\bar d_k$ when $\bar d_k\ge-\hat\omega_k\sqrt{2\ln\ln T/T}$ and at zero otherwise; the upper $p$-value recentres every strategy at its mean.
  \item[StepM.] The stepwise procedure of \citet{romano2005stepwise} on the same studentized statistics and consistent recentring; it returns a set of strategies rather than a $p$-value, so we score whether it declares at least one strategy superior at 5\%.
\end{description}
Finally, an \emph{oracle} knows which strategy is skilled and applies a one-sided Newey--West $t$ test to it alone. No valid best-of-$N$ test can be more powerful, so the oracle's power is a ceiling for every other column.

LET, HC, HC-Lo and DSR call the implementations behind the author's public calculator and API, and RC, SPA and StepM the data-snooping implementation those use; the rest are written in the benchmark. No validator is tuned to the benchmark.

\subsection{Performance measures, Monte Carlo error and runs}
Each cell (family, skill level) of the main run, Null Zoo v1, has \nz{reps.v1} searches. The Monte Carlo standard error of a rejection rate $\hat p$ from $n$ searches is $\sqrt{\hat p(1-\hat p)/n}$: \nz{se5.10000} percentage points at 5\% for $n=10{,}000$. We flag a size as miscalibrated when it lies outside the 99\% band around 5\% (\nz{band99.10000}\% for 10,000 searches), so that a calibrated validator is flagged in about one size cell in a hundred.

The study has three governed runs, each pre-registered before it was run:
\begin{itemize}
  \item \textbf{v1} (base seed \nz{v1.base_seed}). Its pre-registration fixed the families, the skill levels, the number of searches, a size bar (at most \nz{v1.bar_size}\% at a nominal 5\% in every family), a power bar (at least \nz{v1.bar_power}\% at a Sharpe ratio of 2 under i.i.d.\ returns), the oracle ceiling and an ordered list of candidate tests, to choose a default test for the author's open-source validation software. The listed joint-test candidates were \nz{v1.prereg.candidates}. The Reality Check was measured but not listed.
  \item \textbf{v1b} (base seed \nz{v1b.base_seed}). v1 found the Reality Check calibrated where the listed SPA candidates were not. Adopting it on that evidence would have been a choice made after seeing the results, so a separate pre-registration added it to the list and repeated the run on fresh seeds; v1b's result decided the default whichever way it fell.
  \item \textbf{v2} (base seed \nz{v2.base_seed}; commit \texttt{\nz{v2.prereg.commit}} of the repository). An independent implementation in a second language, written from the published definitions rather than translated from v1's code, tests why the studentized SPA test is oversized and where studentizing pays off. Its pre-registration fixed four experiments and \nz{v2.predictions.count} predictions; a pilot run on another seed, which shaped the design, is disclosed in the pre-registration and shares no search with the governed run. Section~\ref{sec:snooping} reports every prediction, including any that failed.
\end{itemize}
v1 and v1b use a xoshiro128** generator \citep{blackman2021scrambled} with one seed per chunk of searches, merged in a fixed order so the output is identical for any number of worker threads; v2 uses numpy's PCG64 generator in the same way. An earlier published version (v0, \nz{reps.v0} searches per cell with seven validators and eight families) and its \nz{reps.ext}-search robustness run are reported in the first version of this paper; v1 supersedes their size and power estimates, and v0 still reproduces byte for byte from its code.
